\section{Exact algorithms for a supplied input}\label{sec:algorithms}

The minimax results quantify the hardest input. We now optimize the schedule
for a supplied input on an arbitrary grid $\mathcal T=(x_0,\ldots,x_N)$.
The data are rational integrated cell allocations
$a_{ji}=A_i(x_j)-A_i(x_{j-1})$, with $a_{ji}\ge0$ and
$\sum_i a_{ji}=\Delta_j$. By Lemma~\ref{lem:endpoint}, these data determine
the error of every grid schedule even when the original relaxed control
varies inside a cell. Arithmetic-operation bounds below count rational
operations; the rational bit lengths are also polynomial in the input length
for each fixed switch budget.

\subsection{One switch and an optimal final block}

Order the modes by nonincreasing terminal mass, breaking ties consistently,
and denote the first two by $h_1,h_2$. Only these leaders, and a third leader
for evaluating omitted maxima, need be selected; a full sort is unnecessary.

\begin{theorem}[Three candidates per boundary]\label{thm:linear-one-switch}
For $n\ge3$ and a fixed internal boundary $t$, an optimal ordered pair of
distinct initial and final modes is among
\[
 (h_1,h_2),\qquad(h_2,h_1),\qquad(p_t,h_1),
 \qquad p_t\in\argmax_{i\notin\{h_1,h_2\}}A_i(t).
\]
For $n=2$ the first two pairs suffice. Comparing these candidates at all
permitted internal boundaries with the best constant schedule solves the
at-most-one-switch problem in $O(nN)$ arithmetic operations. The same bound
holds with a prescribed initial mode and any subset of permitted boundaries.
\end{theorem}
\begin{proof}
By~\eqref{eq:one-switch-error}, the error for $p\to q$ at $t$ is
\[
 \max\{\max_{i\notin\{p,q\}}m_i,\ t-A_p(t),\ T-t-m_q\}.
\]
Replacing $q$ by a largest-mass mode other than $p$ cannot increase any
term: the newly omitted mode has no larger mass than the mode removed from
the omitted set. Thus the final mode can be $h_2$ if $p=h_1$ and $h_1$
otherwise. The first two initial choices give the first two pairs. For
$p\notin\{h_1,h_2\}$ the omitted maximum is $m_{h_2}$ and the final-mode
term is fixed; maximizing $A_p(t)$ gives the third pair.

Select at most three mass leaders in $O(n)$ time and form the cumulative
data in $O(nN)$ operations. Each permitted boundary needs one scan to find
$p_t$ and then a constant number of error evaluations. A largest-mass mode
is the best constant choice, with error $T-m_{h_1}$. If the initial mode is
fixed to $p$, compare only its constant schedule and $p\to q_p$, where
$q_p$ has largest mass among the other modes. With no eligible boundaries,
only the constant comparison remains.
\end{proof}

For a fixed initial mode $p$ and its final choice $q_p$, the two varying
terms are $L_p(t)=t-A_p(t)$ and $R_p(t)=T-t-m_{q_p}$. Their difference
$2t-A_p(t)-T+m_{q_p}$ is strictly increasing, since its increase between
$u<v$ is at least $v-u$. Before its unique crossing, if one exists,
$\max(L_p,R_p)$ decreases; after it, that maximum is nondecreasing.
The omitted-mass term is constant. Consequently the two permitted boundaries
adjacent to the crossing suffice, with the first or last boundary when the
crossing lies outside the permitted range. Binary search over a sorted
boundary list gives $O(n\log(N+1))$ additional queries once the cumulative
data and mass leaders are available for random access. This bound excludes
reading the raw $nN$ input values and preparing the boundary list.

Permitted boundaries can enforce a common minimum duration $d$ for both
nonempty runs by retaining $d\le t\le T-d$, provided constants are allowed
and $T\ge d$. The dominance theorem does not by itself permit arbitrary
mode-specific dwell requirements or restrictions coupling the two labels.
The general algorithm below handles mode-specific minimum dwell times.

\begin{lemma}[Best completion of a fixed prefix]\label{lem:final-residual}
Fix a schedule on $[0,u]$, where $u<T$, with discrepancy $E_0$ and service
$c_i=W_i(u)$. Put $r_i=m_i-c_i$, so that $\sum_i r_i=T-u$.
Completing it in mode $j$ has exact error
\begin{equation}\label{eq:final-residual}
 \max\{E_0,\ \max_{i\ne j}r_i,\ T-u-r_j\}.
\end{equation}
Among any fixed nonempty set of eligible final modes, a mode with largest
residual $r_j$ is optimal. Residuals may be negative.
\end{lemma}
\begin{proof}
Every inactive discrepancy increases to $r_i$, while the active discrepancy
decreases to $r_j-(T-u)$. The other endpoints are already bounded by $E_0$,
which proves the formula. Replacing $j$ by an eligible mode with larger
residual decreases the last term and cannot increase the omitted residual
maximum. This argument is unaffected by negative residuals.
\end{proof}
This lemma optimizes the last block after a fixed prefix; it supplies no
greedy rule for choosing that prefix. The eligible set is held fixed during
this comparison.

\subsection{Assigning modes to a fixed set of blocks}

Fix $k$ positive consecutive blocks with grid endpoint indices
$0=b_0<b_1<\cdots<b_k=N$, and lengths
$\ell_j=x_{b_j}-x_{b_{j-1}}$. For mode $i$ and a subset
$U\subseteq[k]$, define
\begin{equation}\label{eq:block-subset-cost}
 c_i(U)=\max_{1\le j\le k}
 \left|A_i(x_{b_j})-\sum_{\substack{h\in U\\h\le j}}\ell_h\right|.
\end{equation}
This is exactly the component error when mode $i$ occupies those blocks,
including disconnected subsets. In particular $c_i(\varnothing)=m_i$:
unused modes still contribute to the objective.

\begin{proposition}[Subset assignment]\label{prop:subset-DP}
The optimum for these fixed block boundaries is obtained from
\begin{equation}\label{eq:subset-DP}
 D_0(\varnothing)=0,\quad D_0(S)=+\infty\ (S\ne\varnothing),
 \qquad
 D_i(S)=\min_{U\subseteq S}\max\{D_{i-1}(S\setminus U),c_i(U)\},
\end{equation}
as $D_n([k])$. After cumulative preprocessing, its arithmetic cost is
$O(n(k2^k+3^k))$ and its storage, including traceback, is $O(n2^k)$.
\end{proposition}
\begin{proof}
A labeling is a partition of $[k]$ into the subsets assigned to the modes,
with empty subsets allowed. Its discrepancy is the maximum of their costs
by Lemma~\ref{lem:endpoint}. For the first $i$ modes assigning precisely
$S$, enumerate the subset $U$ assigned to mode $i$ and optimize the rest
using the preceding modes. This proves the recurrence by induction.
Each of $2^k$ costs for one mode takes $O(k)$ operations, and
$\sum_{S\subseteq[k]}2^{|S|}=3^k$ counts its transitions. Keep two objective
rows and the minimizing subset for each state and mode. Backtracking from
$D_n([k])$ then gives a labeling with the stated storage.
\end{proof}

\begin{theorem}[Exact fixed-budget optimization]\label{thm:fixed-budget}
Let $s\ge0$ and $k=\min(s+1,N)$. An optimal grid schedule with at most
$s$ switches is computable in
\begin{equation}\label{eq:fixed-budget-complexity}
 O\left(nN+\binom{N-1}{k-1}n(k2^k+3^k)\right)
\end{equation}
arithmetic operations and $O(nN+n2^k)$ storage. In particular, two switches
admit $O(nN^2)$ exact optimization. Repeated modes are permitted.
\end{theorem}
\begin{proof}
Enumerate all $\binom{N-1}{k-1}$ sets of internal block boundaries and use
Proposition~\ref{prop:subset-DP}. Every resulting labeling has at most
$k-1\le s$ changes. Conversely, any feasible grid schedule has at most $k$
maximal runs. Subdivide these runs at unused grid boundaries until there
are exactly $k$ positive nominal blocks, assigning the same label to adjacent
pieces. This is possible because $k\le N$ and does not change the control.
Thus the enumeration covers every feasible schedule. Retain only the best
answer and reuse the workspace for successive partitions.
\end{proof}

For fixed $s$, all rational numbers formed by this algorithm are sums or
differences of polynomially many input numbers; comparisons do not enlarge
them. Their bit lengths remain polynomial. The method is linear in $n$
in its arithmetic count for fixed block count, but its factor of order
$N^{k-1}$ prevents a fixed-parameter tractability claim in $s$ alone.
Switching-time enumeration and exact graph methods for CIA are established
\citep[Section~6.4.3]{zeile2021thesis}\citep{bestehorn2021shortest}.
The subset recurrence is a standard partition dynamic program; its useful
specialization here groups all occurrences of a mode, avoiding enumeration
of the $n^k$ possible mode words.

\subsection{Dwell constraints and a candidate-label bound}

Let each mode have a nonnegative minimum dwell duration $d_i$. We require
every maximal constant run, including initial and terminal runs, to have
length at least $d_i$ when its mode is $i$. An unused mode has no dwell
obligation. In~\eqref{eq:block-subset-cost}, set $c_i(U)=+\infty$ if a
maximal string of consecutive indices in $U$ has total block length below
$d_i$. Consecutive assigned blocks must be merged before applying this test.
The empty subset remains admissible. This check takes $O(k)$ operations,
so both complexity bounds remain valid. The subset recurrence is still
exact because this admissibility depends only on $i$ and $U$. Moreover,
subdivision of a run preserves its physical length and hence its feasibility,
so the proof of Theorem~\ref{thm:fixed-budget} remains valid. The algorithm
detects infeasibility; for example, if every $d_i>T$, no schedule exists.

Other restrictions depending only on one mode's assigned subset can be
treated by the same infinite-cost convention, with their checking cost
included. Extending such a rule from prescribed blocks to the full-budget
enumeration additionally requires invariance under subdivision into adjacent
equally labeled blocks. Arbitrary transition graphs couple different modes
and are outside this separation argument.

An alternative fixes both the block boundaries and their equality pattern.
There are then $d\le k$ distinct roles, each comprising a nonempty subset
of blocks, to be assigned injectively to modes. Let $c_{ri}$ be the
full-horizon component error when role $r$ is assigned to mode $i$;
unselected modes have cost $m_i$.

\begin{lemma}[A bounded candidate set]\label{lem:candidate-labels}
Suppose $n\ge d$. In the unrestricted role-assignment problem, some optimal
assignment uses only the union of the $d$ largest-mass modes and, for each
role $r$, the $d$ modes with smallest $c_{ri}$. This union has at most
$d^2+d$ elements. Ties may be broken consistently in any way.
\end{lemma}
\begin{proof}
Start with an optimal assignment using a label $i$ outside that union. Since
$i$ is not among the $d$ largest-mass labels and there are only $d$ assigned
roles, one of those $d$ largest labels is omitted. Hence $m_i$ is no larger
than the current omitted-mass maximum. Among the $d$ cheapest labels for
$i$'s role, one label $j$ is unused by the other $d-1$ roles. It differs
from $i$ and has no larger role cost. Replace $i$ by $j$. The newly omitted
mass $m_i$ is bounded by the old omitted maximum, so the objective cannot
increase. This removes one label outside the union and adds none. Iteration
proves the claim. If $n=d$, all labels already belong to the union.
\end{proof}
This elementary exchange bound is useful for searches over equality patterns.
The subset algorithm does not need it. No assertion of a new general
assignment principle or unrestricted constraint support is intended.

\subsection{Exact continuous optimization for piecewise-affine input}

For completeness, a finite exact benchmark also exists when switches may
occur inside input cells. Suppose $A$ is piecewise affine on an $N$-cell
rational input grid, and let $k=s+1$ without truncating it at $N$. Fix a word
$p\in[n]^k$ and a nondecreasing list of $k-1$ input-cell indices, one for
each switch time. Write $t_0=0,t_k=T$. On its assigned cell,
$A_i(t_j)$ is affine in $t_j$. The service at that endpoint is
\[
 W_i(t_j)=\sum_{h=1}^j\mathbf1_{\{p_h=i\}}(t_h-t_{h-1}).
\]
Therefore the cell bounds, $t_1\le\cdots\le t_{k-1}$, $0\le E\le T$,
and
\begin{equation}\label{eq:continuous-cell-LP}
 -E\le A_i(t_j)-W_i(t_j)\le E
 \quad(i\in[n],\ 1\le j\le k)
\end{equation}
form a rational linear program in $k$ variables, minimizing $E$.
Lemma~\ref{lem:endpoint} proves that no additional error constraints at
input breakpoints are needed. For the one-sided objective retain only
$W_i(t_j)-A_i(t_j)\le E$.

There are $n^k\binom{N+k-2}{k-1}$ word/cell choices. They cover all schedules:
zero-length blocks and repeated labels are allowed, and a switch at an input
boundary may be assigned to either incident closed cell. The minimum of
their LP values is consequently the continuous instance optimum. Each
feasible LP is bounded and has $O(nk+k)$ inequalities. One entirely rational
solution method enumerates all $k$-element subsets of its inequalities,
solves the corresponding nonsingular equality systems, and retains the best
feasible vertex. A nonempty bounded polytope has an optimal vertex with
$k$ independent active constraints, including when its affine dimension is
smaller than $k$. Thus this method is complete. Its operation count adds a
factor $\binom{O(nk+k)}{k}$ and polynomial work per vertex, and its bit
complexity is polynomial for fixed $k$ by rational elimination and determinant
bounds. This deliberately small-instance benchmark is less efficient in
the number of modes than the grid subset algorithm. Its enumeration principle
is classical; the explicit time-cell decomposition states how rational
certification can be made complete.
